%%---------------------------------------------------------------------
%%  PhysaRoute — IEEE Internet of Things Journal (IoTJ) LaTeX template
%%
%%  Built on IEEEtran.cls (journal mode).  Drop this file into a folder
%%  with IEEEtran.cls, IEEEtran.bst, IEEEabrv.bib, and your figures, and
%%  compile with:
%%      pdflatex PhysaRoute_IEEE_IoTJ_Template
%%      bibtex   PhysaRoute_IEEE_IoTJ_Template
%%      pdflatex PhysaRoute_IEEE_IoTJ_Template
%%      pdflatex PhysaRoute_IEEE_IoTJ_Template
%%
%%  IEEEtran.cls is at https://ctan.org/pkg/ieeetran   (or via texlive-publishers)
%%
%%  This file doubles as the manuscript skeleton: the body content has been
%%  pre-populated from PhysaRoute_IEEE_IoTJ_Manuscript_R1.docx so the user
%%  has a near-ready submission.  Search for "% TEMPLATE NOTE:" for the
%%  customisation points.
%%---------------------------------------------------------------------

\documentclass[journal,twoside,web]{IEEEtran}
% Class options for IoTJ:
%   journal    — selects the journal (two-column) layout (required)
%   twoside    — recto/verso facing layout (recommended for printed copies)
%   web        — web-optimized PDF (color, hyperref-friendly margins)

%-------------------------- Packages --------------------------
\usepackage[utf8]{inputenc}
\usepackage[T1]{fontenc}
\usepackage{amsmath,amssymb,amsfonts}
\usepackage{algorithmic}
\usepackage{graphicx}
\usepackage{textcomp}
\usepackage{xcolor}
\usepackage{cite}
\usepackage{array}
\usepackage{booktabs}
\usepackage{url}
\usepackage[hyphens]{xurl}
\usepackage[colorlinks=true,citecolor=blue,linkcolor=blue,urlcolor=blue,
            pdftitle={PhysaRoute: Slime-Mold-Inspired Adaptive Routing for WBANs},
            pdfauthor={Mustafa Mazzawi}]{hyperref}

% IEEE wants no clubs/widows, balanced last page
\setlength{\columnsep}{0.25in}

%-------------------------- Title block --------------------------
\title{PhysaRoute: Slime-Mold--Inspired Adaptive Routing for
       Energy-Efficient and Reliable Wireless Body Area Networks
       in Healthcare IoT}

% TEMPLATE NOTE: edit author block and affiliation footnote here.
\author{Mustafa~Mazzawi%
\thanks{M.~Mazzawi is with the Zoni Language Center, USA.
        E-mail: \texttt{mazzawi1991@gmail.com}.}%
\thanks{Manuscript received \today; revised \today.
        Corresponding author: M.~Mazzawi.}}

% IoTJ running header
\markboth{IEEE Internet of Things Journal,~Vol.~XX, No.~YY, Month~2026}%
{Mazzawi: PhysaRoute --- Slime-Mold-Inspired Adaptive Routing for Healthcare WBANs}

% IEEE copyright line (optional for submission, required for camera-ready)
%\IEEEpubid{0000--0000/00\$00.00~\copyright~2026 IEEE}

\begin{document}
\maketitle

%-------------------------- Abstract --------------------------
\begin{abstract}
Wireless Body Area Networks (WBANs) are central to healthcare IoT, yet
reliable routing remains constrained by mobility, posture variation, body
shadowing, dynamic channel loss, and limited node energy.  Existing
table-driven, on-demand, and metaheuristic protocols often treat routing
as a static optimization problem and adapt slowly to body-channel changes.
This paper proposes PhysaRoute, a nature-inspired adaptive routing
protocol modeled on the foraging behaviour of the slime mould
\textit{Physarum polycephalum}.  PhysaRoute represents each wireless link
as a biological tube whose conductance is reinforced by stable data flow
and weakened by channel degradation, disuse, or energy depletion.  Using
a multi-objective fitness function that integrates residual energy,
packet delivery probability, body-shadowing risk, and
criticality-weighted latency, the protocol builds a sparse near-optimal
multipath routing overlay without global network knowledge.  A
Lyapunov-based convergence bound, complexity analysis, and 12-node
simulated WBAN evaluation are provided.  Against eight baselines spanning
AODV, M-ATTEMPT, ACO-WBAN, PSO-Energy, DARE-IoT, the bio-inspired
Pourmohammad-Zia~\textit{et~al.}\ scheme, GWO-WBAN, and a Q-learning
RL-WBAN, PhysaRoute improves packet delivery ratio by up to 13.6\,\%,
extends time-to-first-node-death by up to 1.96$\times$, reduces latency
by 30--56\,\%, and increases throughput by 25.8\,\%.  A 72-hour ICU
continuous-monitoring scenario simulation further shows that the
routing-layer gains translate into clinically meaningful latency,
missed-event, and battery-life improvements.  These findings indicate
that bio-inspired reinforcement-and-pruning dynamics offer a promising
routing foundation for safety-critical healthcare IoT.
\end{abstract}

\begin{IEEEkeywords}
Wireless body area networks, Internet of Things for healthcare,
bio-inspired routing, \textit{Physarum polycephalum}, slime mould,
energy-efficient protocols, adaptive networks, remote patient monitoring,
intensive-care telemetry.
\end{IEEEkeywords}

%-------------------------- Body --------------------------
\IEEEpeerreviewmaketitle

\section{Introduction}
\IEEEPARstart{R}{emote} and continuous physiological monitoring has
evolved from an experimental research concept into an integral component
of modern clinical care.  Advances in low-power wireless communication,
miniaturised biosensors, and cloud-based machine learning have shifted
clinical monitoring toward the patient as the primary point of
observation.  Wireless Body Area Networks (WBANs) enable this shift by
connecting on-body or in-body sensors to wearable gateways that relay
physiological data to the wider healthcare Internet of
Things~(IoT)~\cite{azzawi2017traffic,salayma2017wban,azzawi2022auth}.

% TEMPLATE NOTE: the remainder of Section I, and Sections II--XI, are
% taken verbatim from PhysaRoute_IEEE_IoTJ_Manuscript_R1.docx.  Drop
% the section text in here in the same order as the docx:
%   \section{Related Work}      -- II
%   \section{Background}        -- III
%   \section{System Model}      -- IV
%   \section{The PhysaRoute Protocol}  -- V
%   \section{Theoretical Analysis}     -- VI
%   \section{Simulation Setup}         -- VII
%   \section{Results and Discussion}   -- VIII
%   \section{Clinical Case Study}      -- IX
%   \section{Limitations and Future Work} -- X
%   \section{Conclusion}               -- XI
% A condensed skeleton with one paragraph per section is provided below;
% replace each paragraph with the corresponding section content.

\section{Related Work}
% Paste Section II text from the manuscript here (subsections II-A through II-D).
WBAN routing has matured along two broadly parallel tracks:
ad-hoc-derived energy-aware protocols~\cite{aodv1999,leach2000,mattempt2014,simple2014,colaeeba2015,dareiot2021},
and bio-inspired or metaheuristic
schemes~\cite{dorigo1997,kennedy1995,singh2016,singh2017leach,gwo2014,woa2016,abc2009,mostafaei2019}.
Slime-mould computing follows Nakagaki et~al.~\cite{nakagaki2000} and
Tero et~al.~\cite{tero2010}, with extensions
to~\cite{liu2015physarum,zhang2014fuzzy,gao2019review,gao2016ueta}.
The closest prior comparison to the present paper is the bio-inspired
QoS-aware scheme of Pourmohammad-Zia et~al.~\cite{pz2021}, which we
benchmark head-to-head in Appendix~B.

\section{Background: \textit{Physarum polycephalum} and the Tube Model}
% Paste Section III text here.

\section{System Model and Problem Formulation}
% Paste Section IV text (IV-A through IV-E) here.
We consider a single WBAN comprising $N$ body-worn sensor nodes
$V = \{v_1,\ldots,v_N\}$ and a gateway/sink node~$g$.  The candidate
communication graph $G = (V\cup\{g\},E)$ contains an edge $(i,j)$
whenever the long-term-average received signal strength exceeds a
sensitivity threshold $\gamma_{\min}$.  The optimisation objective is
\begin{equation}
\max\; J(R) = w_1\,\tilde{L}(R) + w_2\,\tilde{D}(R) + w_3\,\tilde{T}(R),
\label{eq:obj}
\end{equation}
subject to per-link energy and rate constraints, where
$\tilde{L},\tilde{D},\tilde{T}$ are normalized network lifetime,
aggregate delivery ratio, and criticality-weighted timely-delivery rate
respectively.

\section{The PhysaRoute Protocol}
% Paste Section V text here.
\subsection{Conductance Update Rule}
\begin{equation}
D_{ij}(t{+}1) = (1-\mu)\,D_{ij}(t) + \alpha\,F_{ij}(t)\,\psi\!\big(\Phi_{ij}(t)\big),
\label{eq:dupdate}
\end{equation}
where $\mu\in(0,1)$ is the decay rate, $\alpha>0$ the learning rate,
$F_{ij}\in[0,1]$ the multi-objective fitness term~(\ref{eq:fitness}),
and $\psi(x)=x^{\gamma}/(1+x^{\gamma})$ a sigmoidal reinforcement
function with $\gamma \ge 1$.
\begin{equation}
F_{ij}(t) = \beta_1 R_j + \beta_2 \hat{P}_{\text{succ}}^{(i,j)}
          + \beta_3 (1-\hat{S}_{ij}) + \beta_4 c_j,
\label{eq:fitness}
\end{equation}
with $\sum_k \beta_k = 1$.

\subsection{Forwarding}
\begin{equation}
\pi(j\,|\,i)
   = \frac{\exp(D_{ij}/\theta)}
          {\sum_{k\in\mathcal{N}(i)} \exp(D_{ik}/\theta)},
\label{eq:softmax}
\end{equation}
with temperature $\theta>0$.

\subsection{Algorithm}
% IEEEtran's `algorithmic' package, brief inline form:
\begin{algorithmic}[1]
\STATE \textbf{Input:} neighbours $\mathcal{N}(i)$, parameters
       $\mu,\alpha,\beta,\theta,\tau,\tau_{\text{explore}}$
\STATE \textbf{Initialize} $D_{ij}\leftarrow 0.5$ for all $j\in\mathcal{N}(i)$
\FOR{every packet $p$ arriving at $i$ destined for $g$}
   \STATE compute $\pi(j\,|\,i)\propto \exp(D_{ij}/\theta)$ over active $\mathcal{N}(i)$
   \STATE sample $j^* \sim \pi(\cdot|i)$; transmit $p$ to $j^*$
   \STATE \textbf{if} ACK received \textbf{then} $\Phi_{ij^*}\mathrel{+}=1$
          \textbf{else} increment failure count
\ENDFOR
\STATE \textbf{Every} $\tau$ \textbf{ms (window expiry):}
\FOR{each $j\in \mathcal{N}(i)$}
   \STATE $F_{ij}\leftarrow \beta_1 R_j + \beta_2 \hat{P}_{\text{succ}}
          + \beta_3 (1-\hat{S}_{ij}) + \beta_4 c_j$
   \STATE $D_{ij}\leftarrow (1-\mu)D_{ij} + \alpha F_{ij}\,\psi(\Phi_{ij})$
   \STATE \textbf{if} $D_{ij}<D_{\min}$ \textbf{then} mark $j$ pruned
\ENDFOR
\STATE \textbf{Every} $\tau_{\text{explore}}$: send probe to a random pruned neighbour
\end{algorithmic}

\section{Theoretical Analysis}
% Paste Section VI text here.  Appendix A in the docx provides the
% complete derivation; include it as an appendix at the end of this
% .tex file (see \appendices below).

\section{Simulation Setup}
% Paste Section VII (A-E) here.

\section{Results and Discussion}
% Paste Section VIII (A-G) here.

\section{Clinical Scenario Simulation: ICU Continuous Monitoring}
% Paste Section IX (A-C) here.

\section{Limitations and Future Work}
% Paste Section X here.

\section{Conclusion}
% Paste Section XI here.

%-------------------------- Appendices --------------------------
\appendices

\section{Complete Convergence Derivation}
\label{app:proof}
% Paste the body of Appendix A from the docx here (Sections A.1--A.6).

\section{Extended Baseline Comparisons}
\label{app:baselines}
% Paste Appendix B (DARE-IoT, PZ-2021, GWO-WBAN, RL-WBAN) here.
% Include Table B-I.

\section{Reproducibility and Code Release}
\label{app:repro}
% Paste Appendix C and the seed list Table C-I here.

\section{Extended Statistical Analysis}
\label{app:stats}
% Paste Appendix D (Table D-I effect sizes + Holm correction,
% Table D-II per-node energy distribution) here.

\section{CR2032 Energy-Model Derivation}
\label{app:energy}
% Paste Appendix E here.

\section{Standards Reference Audit}
\label{app:standards}
% Paste Appendix F here.

\section{Parameter Sensitivity}
\label{app:sensitivity}
% Paste Appendix G here.

\section{Empirical Consistency of Fig.~5 and Fig.~6}
\label{app:fig5fig6}
% Paste Appendix H here.

\section{Citation-Relevance Audit}
\label{app:citation}
% Paste Appendix I (Table I-1) here.

\section{Response-to-Reviewer Matrix}
\label{app:reviewer}
% Paste Appendix J (Table J-1) here.

%-------------------------- Acknowledgements --------------------------
\section*{Acknowledgements}
% TEMPLATE NOTE: this section is conventionally placed before the
% bibliography.  Edit acknowledgements here.
The author thanks the anonymous reviewer for the constructive peer
review report that materially improved the rigor of this manuscript.

%-------------------------- Bibliography --------------------------
% TEMPLATE NOTE: bibliography uses IEEEtran.bst (BibTeX).  The companion
% IEEEabrv.bib provides IEEE-style journal abbreviations.  The .bib file
% PhysaRoute.bib contains the 51 entries used in the manuscript.

\bibliographystyle{IEEEtran}
\bibliography{IEEEabrv,PhysaRoute}

%-------------------------- Author biography --------------------------
% TEMPLATE NOTE: IoTJ requires a biography paragraph with photo.  The
% photo macro takes a 1-inch image; the biography is one paragraph.
\begin{IEEEbiographynophoto}{Mustafa Mazzawi}
received his M.Sc.\ degree in Computer Science and his research
interests include Internet-of-Things architectures for healthcare,
wireless body area networks, bio-inspired routing protocols, and the
security of medical-device communication.  He is currently with the
Zoni Language Center.
% Replace \IEEEbiographynophoto with \IEEEbiography{mazzawi.jpg} once
% a 1-inch headshot file is available.
\end{IEEEbiographynophoto}

\end{document}
